\documentclass[10pt,a4paper]{amsart}
\usepackage[margin=19mm]{geometry}
\usepackage{amsmath,amssymb,mathtools}
\usepackage[scr=rsfs,cal=euler]{mathalfa}
\usepackage{xfrac}
\usepackage[unicode,hidelinks,bookmarks=false]{hyperref}
\newcommand{\ie}{i.e.\ }
\newcommand{\cf}{cf.\ }
\newcommand{\resp}{respectively\ }
\newcommand{\Subsection}{\subsection}
\providecommand{\nobreakdash}{\nobreak\hbox{-}}
\setlength{\emergencystretch}{2em}
\multlinegap0pt
\usepackage{bm}
\usepackage{bbm}
\usepackage{dsfont}
\usepackage{xspace}
\usepackage{cancel}
\usepackage{mathtools}
\usepackage{amsthm}
\makeatletter
\@ifundefined{theorem}{\newtheorem{theorem}{Theorem}}{}
\makeatother
\DeclareMathOperator{\Aut}{\mbox{Aut}}
\title[On some semidirect products of skew braces arising in Hopf-G]{On some semidirect products of skew braces arising in Hopf-Galois theory}
\author{Paul J. Truman}
\date{Source: arXiv:2506.04928v1 (CC BY 4.0)}
\begin{document}
\maketitle

\noindent\emph{Source and licence.} On some semidirect products of skew braces arising in Hopf-Galois theory. Version arXiv:2506.04928v1 (CC BY 4.0), distributed under the Creative Commons Attribution 4.0 International licence, as recorded in the arXiv licence metadata for this exact version and shown on the abstract page https://arxiv.org/abs/2506.04928v1. Full rights notice: licenses/arxiv-2506.04928-CC-BY-4.0.txt.\par\smallskip

\noindent\textit{Editorial scope.} Counted unit: the Theorem whose source label is \texttt{thm\_external\_semidirect} in the article (source0.tex lines 219--234, source label \texttt{thm\_external\_semidirect}), delivered together with the article's own complete original proof of it (source0.tex lines 235--274, one \texttt{proof} environment of 40 lines). No mathematics is written, repaired or re-derived: statement and proof are the article's own text line for line. Required context retained uncounted: eqn\_skew\_brace (lines 111--113). Every definition, display and label of those lines is retained unchanged. Omitted from this package and not redistributed: 1--110, 114--218, 275--793. Nothing inside a retained excerpt is omitted or altered: every retained body is delivered exactly as the article's lines are written, so no omission receipt of any kind accompanies this package. Citations: the retained excerpts carry no citation command at all, so no bibliography entry of the article is transcribed and no reference is redistributed. Reference adaptation: none was needed and none was made; every reference inside the delivered unit resolves inside it, so no unresolved reference or citation is produced. Printed slips kept literal: no printed word, symbol, hypothesis or number of the counted statement or of its proof was altered. Layout and class: the article's own class is \texttt{amsart} with packages including amsmath, amssymb, amsthm, eufrak, setspace, verbatim, xypic, enumitem, geometry, xy, lineno; the wrapper is the lane's pinned \texttt{amsart} editorial wrapper at 10pt on A4, which restates the theorem-like environments the retained text needs and adds the article's own macro definitions that the retained text needs (\texttt{\textbackslash Aut}), with extra wrapper packages (bm, bbm, dsfont, xspace, cancel, mathtools). The delivered numbering is the unit's own. Assets: no retained span contains a figure environment, a graphics-inclusion command, a PDF-page inclusion, a table or a TikZ picture; no image file is copied into this package, the package declares no asset dependency and no image file is redistributed with it. Licence: version arXiv:2506.04928v1 of this article is distributed under the Creative Commons Attribution 4.0 International licence, as recorded in the arXiv licence metadata for this exact version and shown on the abstract page https://arxiv.org/abs/2506.04928v1. A full rights notice is delivered at licenses/arxiv-2506.04928-CC-BY-4.0.txt. Characters: every retained excerpt is pure ASCII, so the delivered engine, XeLaTeX with its default Latin Modern text font, reports no missing character.
\begin{equation} \label{eqn_skew_brace}
x \circ (y \cdot z) = (x \circ y) \cdot x^{-1} \cdot (x \circ z) \mbox{ for all } x,y,z \in G, 
\end{equation}

\begin{theorem} \label{thm_external_semidirect}
Let $ A=(A,\cdot,\circ) $ and $ B=(B,\cdot,\circ) $ be skew braces and let
\begin{eqnarray*}
\varphi : (B,\circ) \rightarrow \Aut(A,\circ) \\
\theta : (B,\cdot) \rightarrow \Aut(A,\cdot)
\end{eqnarray*}
be group homomorphisms. The following statements are equivalent:
\begin{enumerate}
\item \label{enum_external_sdp_1} $ A \rtimes_{\theta}^{\varphi} B $ is a skew brace;
\item \label{enum_external_sdp_2} The skew brace $ B $ acts on the skew brace $ A $ via $ \varphi $ and the following equations hold inside $ \Aut(A,\cdot) $: 
\begin{eqnarray}
\gamma_{a} \theta_{b} = \theta_{b} \gamma_{a} \mbox{ for all } a \in A \mbox{ and } b \in B; \label{eqn_external_sdp_1} \\
\varphi_{b} \theta_{b'} = \theta_{b \gamma_{b}(b') b^{-1}} \varphi_{b} \mbox{ for all } b, b' \in B. \label{eqn_external_sdp_2}
\end{eqnarray}
\end{enumerate}
\end{theorem}

\begin{proof}
First suppose that \ref{enum_external_sdp_1} holds. Then we have
\begin{eqnarray*}
(e,b) \circ (a,e)(a',e) & = & (e,b) \circ (aa',e) \\
& = & (\varphi_{b}(aa'),b), 
\end{eqnarray*}
but also
\begin{eqnarray*}
(e,b) \circ (a,e)(a',e) & = & [(e,b) \circ (a,e)] (e,b)^{-1} [(e,b) \circ (a',e)] \mbox{ by \eqref{eqn_skew_brace} } \\
& = & (\varphi_{b}(a),b)(e,b^{-1}) (\varphi_{b}(a'),b) \\
& = & (\varphi_{b}(a)\varphi_{b}(a'),b).
\end{eqnarray*}
Hence $ \varphi_{b}(aa') = \varphi_{b}(a)\varphi_{b}(a') $, so $ \varphi_{b} \in \Aut(A,\cdot) $, and so the skew brace $ B $ acts on the skew brace $ A $. 

To complete the proof we assume that $ B $ acts on $ A $ via $ \varphi $ and show that the skew brace relation is satisfied on $ A \rtimes_{\theta}^{\varphi} B $ if and only if \eqref{eqn_external_sdp_1} and \eqref{eqn_external_sdp_2} hold. Let $ a,c,c' \in A $ and $ b,d,d' \in B $. Then we have
\begin{align}
(a,b) \circ (c,d)(c',d') &= (a,b) \circ (c \theta_{d}(c'), dd') \nonumber \\
&= (a \circ \varphi_{b}(c\theta_{d}(c')), b \circ dd') \nonumber \\
&= (a \circ \varphi_{b}(c)\varphi_{b}(\theta_{d}(c')), b \circ dd') \tag{ since   $\varphi_{b} \in \Aut(A,\cdot)$}\nonumber \\
&= ((a \circ \varphi_{b}(c)) a^{-1} (a \circ (\varphi_{b}(\theta_{d}(c'))), b \circ dd') \nonumber \\
&= ((a \circ \varphi_{b}(c)) \gamma_{a}(\varphi_{b}(\theta_{d}(c'))), b \circ dd') \label{eqn_external_sdp_proof_1}
\end{align}
whereas
\begin{eqnarray}
[(a,b) \circ (c,d)](a,b)^{-1}[(a,b) \circ (c',d')] 
& = & (a \circ \varphi_{b}(c), b \circ d)(\theta_{b}^{-1}(a^{-1}),b^{-1})(a \circ \varphi_{b}(c'),b \circ d') \nonumber \\
& = & (a \circ \varphi_{b}(c), b \circ d)(\theta_{b}^{-1}(a^{-1})\theta_{b}^{-1}(a \circ \varphi_{b}(c')),b^{-1}(b \circ d')) \nonumber \\
& = & (a \circ \varphi_{b}(c), b \circ d)(\theta_{b}^{-1}(a^{-1}(a \circ \varphi_{b}(c')),b^{-1}(b \circ d')) \nonumber \\
& = & (a \circ \varphi_{b}(c), b \circ d)(\theta_{b}^{-1}(\gamma_{a}(\varphi_{b}(c')),b^{-1}(b \circ d')) \nonumber \\
& = & ((a \circ \varphi_{b}(c)) \theta_{b \circ d}\theta_{b}^{-1}(\gamma_{a}(\varphi_{b}(c')),(b \circ d)b^{-1}(b \circ d')) \nonumber \\
& = & ((a \circ \varphi_{b}(c)) \theta_{b}\theta_{\gamma_{b}(d)}\theta_{b}^{-1}(\gamma_{a}(\varphi_{b}(c')),b \circ dd') \label{eqn_external_sdp_proof_2}
\end{eqnarray}
Hence $ A \rtimes_{\theta}^{\varphi} B $ is a skew brace if and only if \eqref{eqn_external_sdp_proof_1} and \eqref{eqn_external_sdp_proof_2} agree, which occurs if and only if
\begin{equation} \label{eqn_external_sdp_proof_3}
\gamma_{a}\varphi_{b}\theta_{d} = \theta_{b\gamma_{b}(d)b^{-1}}\gamma_{a}\varphi_{b} \in \Aut(A) \mbox{ for all } a \in A, b \in B.
\end{equation}
If \eqref{eqn_external_sdp_proof_3} holds then choosing $ b=1 $ yields $ \gamma_{a} \theta_{d} = \theta_{d}\gamma_{a} $, and so \eqref{eqn_external_sdp_1} holds; similarly choosing $ a=1 $ yields $ \varphi_{b}\theta_{d} = \theta_{b\gamma_{b}(d)b^{-1}}\varphi_{b} $, so \eqref{eqn_external_sdp_2} holds. Conversely, if \eqref{eqn_external_sdp_1} and \eqref{eqn_external_sdp_2} hold then \eqref{eqn_external_sdp_proof_3} holds. 

This completes the proof that \ref{enum_external_sdp_1} and \ref{enum_external_sdp_2} are equivalent.
\end{proof}

\end{document}
